\documentclass[11pt]{article}
\usepackage[margin=1in]{geometry}
\usepackage{amsmath}
\usepackage{amssymb}
\usepackage{booktabs}
\usepackage{longtable}
\usepackage{hyperref}
\usepackage{microtype}

\title{Fluo-suite nuclei-selection survey: connected-nuclei counts vs.\ dipolar cutoff}
\author{}
\date{}

\begin{document}
\maketitle

\section{Method}

For each molecule in the fluo\_suite dataset
(\url{circ_sim/data/fluo_suite/ORCA_NMR_summary_merged.xlsx}),
a chemistry-constrained candidate pool is built from the full set of $^{1}$H/$^{13}$C/$^{19}$F
nuclei as follows (see \url{circ_sim/scripts/zulf_numerics/fluo_suite/moment_convergence.py}):

\begin{enumerate}
\item \textbf{Seed}: all $^{19}$F atoms, plus a single representative carbon,
  reflecting natural $^{13}$C abundance ($\sim$1.1\,\%): at most one carbon
  is isotopically active in a given molecule at a time.
\item \textbf{Candidate pool}: all non-exchangeable protons, identified as
  H atoms with at least one 1-bond C--H coupling in the dataset (a
  data-driven proxy for ``bonded to carbon, not to O/N''). Exchangeable
  protons (–OH, –NH$_2$, \ldots) are excluded.
\end{enumerate}

\paragraph{Multiple fluorinated anchors.} A ``fluorinated anchor'' is a
carbon directly (1-bond) coupled to at least one F atom. Several molecules
in the suite contain more than one \emph{chemically separate} anchor (e.g.\
two independent CF$_2$ or CF$_3$ groups on different rings). Since natural
abundance only ever makes one carbon isotopically active at a time, each
anchor is an equally valid, mutually exclusive scenario for ``which carbon
is the $^{13}$C this time'' -- there is no principled reason to keep only
the most strongly $J$-coupled one. When a molecule has $N\geq 1$ such
anchors, $N$ separate versions of it are therefore carried through the rest
of the analysis, one per anchor, labelled with a letter suffix (a, b, c,
\ldots) ordered by decreasing cumulative $J(\mathrm{C,F})$ -- version ``a''
reproduces the single-anchor choice used in the original (pre-versioning)
survey. In every version the seed still contains \emph{all} $^{19}$F atoms
in the molecule, not just the ones bonded to that version's anchor: an
F-detected experiment observes every $^{19}$F resonance simultaneously
regardless of which carbon happens to be $^{13}$C, so only the local
splitting/connectivity pattern near the labelled anchor differs between
versions, not the detection channel itself.

The \emph{maximal number of chemically/NMR-relevant nuclei} for a molecule
is the size of this constrained set (seed + candidate pool), independent of
any distance cutoff -- the ceiling on how many nuclei could ever be
retained under these two chemistry constraints alone.

Within that constrained set, reachability from the seed is then determined
by a combined coherent+dipolar coupling graph:
\begin{itemize}
\item \emph{Coherent edges}: 1-bond/2-bond $J$-couplings from the ORCA/PySCF
  dataset, weight $2\pi|J_{ij}|$ (rad/s).
\item \emph{Dipolar edges}: through-space coupling computed from atomic
  coordinates for \emph{every} pair within the stated cutoff, weight
  $b_{ij}^{2}\cdot(2\tau_c/5)$ (rad/s), with
  $b_{ij}=-(\mu_0/4\pi)\gamma_i\gamma_j\hbar/r_{ij}^{3}$ and
  $\tau_c=10^{-10}\,\mathrm{s}$ -- the same spectral-density prefactor used
  in the rank-2 Lindblad jump operators elsewhere in this codebase, so the
  two channels are combined in matched units rather than an arbitrary
  rescaling.
\end{itemize}
A candidate atom is counted as \emph{connected} if there exists a path of
nonzero-weight edges (either channel) back to the seed. This reachability
criterion is a graph-only upper bound: it does not verify that a reachable
atom measurably perturbs the $^{19}$F-detected signal (see the
moment-convergence discussion in \texttt{moment\_convergence.py} for a
tighter, dynamics-based criterion).

\section{Results}

Table~\ref{tab:connectivity} reports, for every molecule in the suite
(mol4, Pibrentasvir, was not computed at the ORCA/PySCF stage and is
omitted) and every fluorinated-anchor version thereof, the anchor carbon
(and the F atom(s) it is 1-bond coupled to), the total number of NMR-active
nuclei, the maximal chemistry-relevant set size, and the number of nuclei
connected to the seed under the combined graph at cutoffs of 5.0, 3.0, 2.5,
and 2.0\,\AA. 12 of the 21 molecules have more than one fluorinated anchor
(up to 4, for Atogepant), giving 36 rows in total.

\small
\setlength{\tabcolsep}{3pt}
\begin{longtable}{@{}l l l r r r r r r@{}}
\toprule
ID & Drug name & Anchor & Total nuclei & Max relevant & \multicolumn{4}{c}{Connected nuclei} \\
   &           &        & (H+C+F)      & (ceiling)    & 5.0\,\AA & 3.0\,\AA & 2.5\,\AA & 2.0\,\AA \\
\midrule
\endhead
mol1     & Lorlatinib               & C6 (F8)                &  41 &  19 &  19 &  15 &  13 &   4 \\
mol2a    & Voxilaprevir             & C40 (F54,F55)          &  96 &  54 &  54 &  48 &  21 &   8 \\
mol2b    &                          & C13 (F14,F15)          &  96 &  54 &  54 &  48 &  21 &   7 \\
mol3a    & Glecaprevir              & C18 (F19,F20)          &  88 &  48 &  48 &  44 &  26 &   7 \\
mol3b    &                          & C37 (F38,F39)          &  88 &  48 &  48 &  44 &  30 &   7 \\
mol5a    & Belzutifan               & C16 (F22)              &  32 &  15 &  15 &   9 &   9 &   8 \\
mol5b    &                          & C8 (F23)               &  32 &  15 &  15 &   9 &   7 &   6 \\
mol5c    &                          & C7 (F24)               &  32 &  15 &  15 &   9 &   7 &   7 \\
mol6     & Pexidartinib             & C25 (F26,F27,F28)      &  38 &  17 &  17 &   9 &   5 &   4 \\
mol7     & Tecovirimat              & C23 (F24,F25,F26)      &  37 &  18 &   8 &   8 &   5 &   4 \\
mol8a    & Ivosidenib               & C29 (F31,F32)          &  53 &  25 &  25 &  23 &   8 &   8 \\
mol8b    &                          & C11 (F15)              &  53 &  25 &  25 &  23 &   6 &   6 \\
mol9a    & Tezacaftor               & C33 (F35,F36)          &  56 &  27 &  27 &  19 &   4 &   4 \\
mol9b    &                          & C10 (F19)              &  56 &  27 &  27 &  19 &  19 &   5 \\
mol10    & Emtricitabine            & C9 (F15)               &  19 &   9 &   9 &   9 &   9 &   3 \\
mol11    & Gemcitabine              & C9 (F16,F17)           &  22 &  10 &  10 &  10 &   7 &   5 \\
mol12    & Sofosbuvir               & C15 (F27)              &  52 &  28 &  28 &  21 &   9 &   7 \\
mol13    & Tedizolid phosphate      & C13 (F30)              &  34 &  16 &  13 &   9 &   9 &   3 \\
mol14a   & Cabotegravir             & C23 (F26)              &  38 &  18 &  18 &   8 &   5 &   5 \\
mol14b   &                          & C21 (F27)              &  38 &  18 &  18 &   8 &   4 &   4 \\
mol15    & Linezolid                & C13 (F23)              &  37 &  21 &  21 &  18 &  18 &   3 \\
mol16    & Eravacycline             & C9 (F39)               &  59 &  26 &  26 &  16 &   9 &   3 \\
mol17a   & Atogepant                & C9 (F10,F11,F12)       &  58 &  28 &  28 &  19 &  17 &   9 \\
mol17b   &                          & C38 (F41)              &  58 &  28 &  28 &  19 &  18 &   8 \\
mol17c   &                          & C39 (F40)              &  58 &  28 &  28 &  19 &  17 &   7 \\
mol17d   &                          & C35 (F42)              &  58 &  28 &  28 &  19 &  18 &   8 \\
mol18a   & Letermovir               & C8 (F9,F10,F11)        &  61 &  32 &  32 &  32 &  10 &   5 \\
mol18b   &                          & C16 (F20)              &  61 &  32 &  32 &  32 &  11 &   6 \\
mol19    & Siponimod                & C24 (F25,F26,F27)      &  67 &  38 &  38 &  20 &  20 &   4 \\
mol20a   & Fluticasone furoate      & C22 (F23)              &  59 &  32 &  32 &  32 &  28 &   7 \\
mol20b   &                          & C14 (F34)              &  59 &  32 &  32 &  32 &  28 &   9 \\
mol20c   &                          & C6 (F36)               &  59 &  32 &  32 &  32 &  28 &   9 \\
mol21a   & Fluocinolone acetonide   & C4 (F30)               &  56 &  31 &  31 &  31 &  31 &   6 \\
mol21b   &                          & C20 (F29)              &  56 &  31 &  31 &  31 &  31 &   6 \\
mol22a   & Dutasteride              & C19 (F20,F21,F22)      &  63 &  35 &  35 &  35 &  35 &   7 \\
mol22b   &                          & C23 (F24,F25,F26)      &  63 &  35 &  35 &  35 &  35 &   7 \\
\bottomrule
\caption{Connected-nuclei counts under the combined coherent+dipolar graph,
by cutoff, for every fluorinated-anchor version of every molecule in the
suite. ``Anchor'' gives the version's 13C atom and the F atom(s) it is
1-bond coupled to. ``Max relevant'' is the chemistry-constrained ceiling
(seed + non-exchangeable-H candidate pool) and, like ``Total nuclei'', is
identical across versions of the same molecule (it depends only on the
number of F atoms and non-exchangeable H's, not on which carbon is chosen).
At 5.0\,\AA\ the graph reaches the ceiling for most molecule/anchor
versions; mol7 (Tecovirimat) and mol13 (Tedizolid phosphate) are exceptions
where part of the candidate pool sits 10+\,\AA\ from the seed even at this
cutoff (an elongated molecule with the anchor at one end and remaining
protons on a distant ring), not a data or algorithm artifact. By 2.5\,\AA\
the different anchors of the same molecule already start to diverge
noticeably (e.g.\ mol9a: 4 vs.\ mol9b: 19; mol3a: 26 vs.\ mol3b: 30), well
before either reaches its 2.0\,\AA\ value -- so 2.5\,\AA\ is where
anchor-dependence first becomes significant, not just a smaller version of
the 5.0\,\AA\ picture. At 2.0\,\AA\ almost all versions collapse to a small
(3--9-nuclei) immediate neighborhood of the anchor, where different anchors
of the same molecule continue to disagree (e.g.\ mol5: 8/6/7, mol17:
9/8/7/8, mol20: 7/9/9) -- consistent with tight cutoffs only reaching each
anchor's own local environment.\label{tab:connectivity}}
\end{longtable}

\section{Filtered cutoff selection}

Table~\ref{tab:filtered} applies a simple rule to pick, per molecule/anchor
version, a small pair of cutoffs to actually use rather than the full
5.0/3.0/2.5/2.0\,\AA\ sweep: always keep 5.0\,\AA, then take the \emph{next}
cutoff (scanning 3.0\,\AA, then 2.5\,\AA, then 2.0\,\AA) at which the
connected-nuclei count first changes from its 5.0\,\AA\ value. If that
changed value is $\geq 9$ it is kept as the second cutoff; if it is $<9$,
no second cutoff is kept for that row (only 5.0\,\AA\ is reported) --
scanning does not continue past the first change point even if a later,
untested-in-this-rule cutoff might happen to be $\geq 9$, since the count is
monotonically non-increasing as the cutoff tightens, so nothing later in
the scan could exceed a value already below the floor. Rows whose
5.0\,\AA\ count itself is $<9$ are dropped entirely, since no cutoff in the
sweep could satisfy the floor -- this removes mol7 (Tecovirimat: 8 at
5.0\,\AA\ already), leaving 35 of the 36 rows.

\small
\setlength{\tabcolsep}{3pt}
\begin{longtable}{@{}l l l r r r r@{}}
\toprule
ID & Drug name & Anchor & Total nuclei & Connected & Second & Connected \\
   &           &        & (H+C+F)      & @5.0\,\AA & cutoff (\AA) & @second cutoff \\
\midrule
\endhead
mol1     & Lorlatinib               & C6 (F8)                &  41 &  19 & 3.0  &  15 \\
mol2a    & Voxilaprevir             & C40 (F54,F55)          &  96 &  54 & 3.0  &  48 \\
mol2b    &                          & C13 (F14,F15)          &  96 &  54 & 3.0  &  48 \\
mol3a    & Glecaprevir              & C18 (F19,F20)          &  88 &  48 & 3.0  &  44 \\
mol3b    &                          & C37 (F38,F39)          &  88 &  48 & 3.0  &  44 \\
mol5a    & Belzutifan               & C16 (F22)              &  32 &  15 & 3.0  &   9 \\
mol5b    &                          & C8 (F23)               &  32 &  15 & 3.0  &   9 \\
mol5c    &                          & C7 (F24)               &  32 &  15 & 3.0  &   9 \\
mol6     & Pexidartinib             & C25 (F26,F27,F28)      &  38 &  17 & 3.0  &   9 \\
mol8a    & Ivosidenib               & C29 (F31,F32)          &  53 &  25 & 3.0  &  23 \\
mol8b    &                          & C11 (F15)              &  53 &  25 & 3.0  &  23 \\
mol9a    & Tezacaftor               & C33 (F35,F36)          &  56 &  27 & 3.0  &  19 \\
mol9b    &                          & C10 (F19)              &  56 &  27 & 3.0  &  19 \\
mol10    & Emtricitabine            & C9 (F15)               &  19 &   9 & --   &  -- \\
mol11    & Gemcitabine              & C9 (F16,F17)           &  22 &  10 & --   &  -- \\
mol12    & Sofosbuvir               & C15 (F27)              &  52 &  28 & 3.0  &  21 \\
mol13    & Tedizolid phosphate      & C13 (F30)              &  34 &  13 & 3.0  &   9 \\
mol14a   & Cabotegravir             & C23 (F26)              &  38 &  18 & --   &  -- \\
mol14b   &                          & C21 (F27)              &  38 &  18 & --   &  -- \\
mol15    & Linezolid                & C13 (F23)              &  37 &  21 & 3.0  &  18 \\
mol16    & Eravacycline             & C9 (F39)               &  59 &  26 & 3.0  &  16 \\
mol17a   & Atogepant                & C9 (F10,F11,F12)       &  58 &  28 & 3.0  &  19 \\
mol17b   &                          & C38 (F41)              &  58 &  28 & 3.0  &  19 \\
mol17c   &                          & C39 (F40)              &  58 &  28 & 3.0  &  19 \\
mol17d   &                          & C35 (F42)              &  58 &  28 & 3.0  &  19 \\
mol18a   & Letermovir               & C8 (F9,F10,F11)        &  61 &  32 & 2.5  &  10 \\
mol18b   &                          & C16 (F20)              &  61 &  32 & 2.5  &  11 \\
mol19    & Siponimod                & C24 (F25,F26,F27)      &  67 &  38 & 3.0  &  20 \\
mol20a   & Fluticasone furoate      & C22 (F23)              &  59 &  32 & 2.5  &  28 \\
mol20b   &                          & C14 (F34)              &  59 &  32 & 2.5  &  28 \\
mol20c   &                          & C6 (F36)               &  59 &  32 & 2.5  &  28 \\
mol21a   & Fluocinolone acetonide   & C4 (F30)               &  56 &  31 & --   &  -- \\
mol21b   &                          & C20 (F29)              &  56 &  31 & --   &  -- \\
mol22a   & Dutasteride              & C19 (F20,F21,F22)      &  63 &  35 & --   &  -- \\
mol22b   &                          & C23 (F24,F25,F26)      &  63 &  35 & --   &  -- \\
\bottomrule
\caption{Filtered cutoff selection: 5.0\,\AA\ plus the next cutoff at which
the connected-nuclei count changes, subject to a floor of 9 nuclei (see
text). ``--'' means no second cutoff qualifies: either the count never
changes down to 2.0\,\AA\ (mol10, mol11, mol14a/b), or the first change
already drops to or below the floor (mol21a/b at 2.0\,\AA\ goes to 6;
mol22a/b to 7). For most rows the qualifying second cutoff is 3.0\,\AA;
mol18a/b and mol20a/b/c are the only rows where the count survives
unchanged through 3.0\,\AA\ and only breaks at
2.5\,\AA.\label{tab:filtered}}
\end{longtable}

\end{document}
